\subsection{阿廷代数的分裂}

在本小节中，K 是一个交换环，A 是一个 K-代数。用 r 表示 A 的根。我们把商代数 \(A/r\) 记作 \(\overline{A}\)，把从 A 到 \(\overline{A}\) 的典范映射记作 \(\pi\)。我们关注的是 A 的满足 \(A = S \oplus r\) 的子代数 S。

我们用 \(\Sigma\) 表示 K-线性截面 \(s\)（\(\pi\) 的）中满足 \(s(\alpha\beta) = s(\alpha)s(\beta)\)（对 \(\alpha, \beta\) 属于 \(\overline{\mathrm{A}}\)）者所成之集。注意这样的截面必然满足 \(s(1) = 1\)（换言之，\(s\) 是一个环同态）：我们确有 \(s(1)^2 = s(1)\)，且 \(s(1)\) 是可逆的，因为它属于 \(1 + \mathfrak{r}\)（VIII, 第 156 页，定理 1）。若 \(s\) 是 \(\Sigma\) 的一个元素，则 \(s\) 的像 S 是 A 的一个子代数，且我们有 \(\mathrm{A} = \mathrm{S} \oplus \mathfrak{r}\)。反之，若 S 是 A 的一个满足 \(\mathrm{A} = \mathrm{S} \oplus \mathfrak{r}\) 的子代数，则 \(\pi\) 在 S 上的限制是双射的，且该逆双射定义了 \(\Sigma\) 的一个以 S 为像的元素。

由雅各布森定理（同前所引），\(1 + \mathfrak{r}\) 的每个元素在 A 中都是可逆的。我们称形如 \(a \mapsto xax^{-1}\)（其中 \(x \in 1 + \mathfrak{r}\)）的 A 的内自同构为特殊自同构。

命题 7. —— 假设 \((\overline{\mathrm{A}}\otimes_{\mathrm{K}}\overline{\mathrm{A}}^{\mathrm{o}})\) -模 \(\overline{\mathrm{A}}\) 是投射的。

\begin{enumerate}
\def\labelenumi{\alph{enumi})}
\item
  设 \(S_1\) 与 \(S_2\) 是 A 的满足 \(A = S_1 \oplus \mathfrak{r} = S_2 \oplus \mathfrak{r}\) 的子代数。存在 A 的一个把 \(S_1\) 变到 \(S_2\) 的特殊自同构。
\item
  假设 \(\pi\) 有一个 K-线性截面，且 A 的根 r 是幂零的。则存在 A 的一个满足 \(A = S \oplus r\) 的子代数 S。
\end{enumerate}

设 \(S_{1}\) 与 \(S_{2}\) 如 a) 中所述。设 \(s_{1}\) 与 \(s_{2}\) 是集合 \(\Sigma\) 中对应于子代数 \(S_{1}\) 与 \(S_{2}\) 的元素。设 \(\varepsilon\) 是从 \(A \otimes_{K} A\) 到 \(A\) 的由 \(\varepsilon(a \otimes b) = ab\) 给出的 K-线性映射。由 VIII, 第 236 页的命题 5 与 VIII, 第 237 页的注 1，存在元素 \(e = \sum_{i=1}^{r} \alpha_{i} \otimes \alpha_{i}'\)（属于 \(\overline{A} \otimes_{K} \overline{A}\)），满足 \(\sum_{i=1}^{r} \alpha_{i} \alpha_{i}' = 1\) 且 \(\sum_{i=1}^{r} \alpha \alpha_{i} \otimes \alpha_{i}' = \sum_{i=1}^{r} \alpha_{i} \otimes \alpha_{i}' \alpha\)（对每个 \(\alpha \in \overline{A}\)）。令 \(x = \sum_{i=1}^{r} s_{1}(\alpha_{i}) s_{2}(\alpha_{i}')\)。我们有 \(\pi(x) = \sum_{i=1}^{r} \alpha_{i} \alpha_{i}' = 1\)，因此 \(x \in 1 + r\)。设 \(\alpha\) 是 \(\overline{A}\) 的一个元素。我们有

\[
\begin{array}{l} s _ {1} (\alpha) x = \sum_ {i = 1} ^ {r} s _ {1} (\alpha \alpha_ {i}) s _ {2} (\alpha_ {i} ^ {\prime}) = (\varepsilon \circ (s _ {1} \otimes s _ {2})) \biggl (\sum_ {i = 1} ^ {r} \alpha \alpha_ {i} \otimes \alpha_ {i} ^ {\prime} \biggr) \\ = (\varepsilon \circ (s _ {1} \otimes s _ {2})) \biggl (\sum_ {i = 1} ^ {r} \alpha_ {i} \otimes \alpha_ {i} ^ {\prime} \alpha \biggr) = \sum_ {i = 1} ^ {r} s _ {1} (\alpha_ {i}) s _ {2} (\alpha_ {i} ^ {\prime} \alpha) = x s _ {2} (\alpha). \end{array}
\]

由此得到等式 \(x^{-1}S_{1}x = S_{2}\)，从而得到断言 a)。

在 b) 的假设下，再假设 \(r^{2}=0\)。在这种情形，(A, A)-双模 r 被 r 零化，因此我们把它看作一个 \((\overline{A},\overline{A})\) -双模。选取一个 K-线性截面 \(\sigma\)（\(\pi\) 的）。我们有

\[
\alpha x = \sigma (\alpha) x \quad \text { and } \quad x \alpha = x \sigma (\alpha)\tag{19}
\]

对 \(\alpha \in \overline{\mathrm{A}}\) 与 \(x\in \mathfrak{r}\) 成立。令

\[
\varphi (\alpha , \beta) = \sigma (\alpha \beta) - \sigma (\alpha) \sigma (\beta)\tag{20}
\]

对 \(\alpha, \beta \in \overline{\mathbf{A}}\) 成立。我们有关系式 \(\pi(\varphi(\alpha, \beta)) = \alpha\beta - \alpha\beta = 0\)（对 \(\alpha, \beta \in \overline{\mathbf{A}}\)）。因此，\(\varphi\) 定义了 \(\mathrm{C}^2(\overline{\mathbf{A}}, \mathfrak{r})\) 的一个元素。设 \(\alpha, \beta, \gamma\) 是 \(\overline{\mathbf{A}}\) 的元素；注意到 (19)，我们有

\[
\begin{array}{r l} \partial^ {2} \varphi (\alpha , \beta , \gamma) & = \alpha \varphi (\beta , \gamma) - \varphi (\alpha \beta , \gamma) + \varphi (\alpha , \beta \gamma) - \varphi (\alpha , \beta) \gamma \\ & = \sigma (\alpha) \varphi (\beta , \gamma) - \varphi (\alpha \beta , \gamma) + \varphi (\alpha , \beta \gamma) - \varphi (\alpha , \beta) \sigma (\gamma) \\ & = \sigma (\alpha) (\sigma (\beta \gamma) - \sigma (\beta) \sigma (\gamma)) - \sigma (\alpha \beta \gamma) + \sigma (\alpha \beta) \sigma (\gamma) + \sigma (\alpha \beta \gamma) \\ & - \sigma (\alpha) \sigma (\beta \gamma) - (\sigma (\alpha \beta) - \sigma (\alpha) \sigma (\beta)) \sigma (\gamma) \\ & = 0. \end{array}
\]

由 VIII, 第 242 页的定理 3，K-模 \(\mathrm{H}^{2}(\overline{\mathrm{A}},\mathfrak{r})\) 化为零。因此存在元素 \(\psi\)（\(\mathrm{C}^{1}(\overline{\mathrm{A}},\mathfrak{r})\) 的），使得 \(\partial^{1}\psi = \varphi\)，换句话说，使得我们有

\[
\varphi (\alpha , \beta) = \alpha \psi (\beta) - \psi (\alpha \beta) + \psi (\alpha) \beta \quad \text { for } \alpha , \beta \text { in } \overline {{\mathrm{A}}}.\tag{21}
\]

我们有 \(\psi(\alpha)\psi(\beta)=0\)，因为 \(r^{2}\) 为零；于是由 (19) 与 (20) 我们推出

\[
(\sigma + \psi) (\alpha \beta) = (\sigma + \psi) (\alpha) (\sigma + \psi) (\beta),\tag{22}
\]

于是 K-线性截面 \(\sigma +\psi\)（\(\pi\) 的）属于 \(\Sigma\)。它的像是一个子代数 \(S\)（\(A\) 的），满足 \(A = S + \mathfrak{r}\)。

我们现在来证明一般情形中 S 的存在性。我们对最小整数 \(p \geqslant 1\)（使 \(r^{p} = 0\)）作归纳推理；p = 1 的情形是平凡的。假设 \(p \geqslant 2\)，并令 \(A' = A/r^{p-1}\)。记 \(r'\) 为 \(A'\) 的根，它等于 \(r/r^{p-1}\)（VIII, 第 155 页的命题 5），故满足 \(r'^{p-1} = 0\)，且代数 \(A'/r'\) 同构于 \(\overline{A} = A/r\)，从而是绝对半单的。由归纳假设，存在子代数 \(S'\)（\(A'\) 的），使得 \(A' = S' \oplus r'\)。于是 \(S'\) 具有形式 \(A''/r^{p-1}\)，其中 \(A''\) 是 A 的一个包含 \(r^{p-1}\) 的子代数，且我们有

\[
\mathrm{A} = \mathrm{A} ^ {\prime \prime} + \mathfrak {r}, \quad \mathfrak {r} ^ {p - 1} = \mathrm{A} ^ {\prime \prime} \cap \mathfrak {r}.\tag{23}
\]

代数 \(A^{\prime\prime}/r^{p-1}\) 同构于 \(A^{\prime}/r^{\prime}\)；我们有 \((\mathfrak{r}^{p-1})^{2}=0\)，故 \(r^{p-1}\) 是 \(A^{\prime\prime}\) 的根。由刚才处理过的情形，存在 \(A^{\prime\prime}\) 的一个子代数 S，使得 \(A^{\prime\prime}=S\oplus r^{p-1}\)；由 (23) 我们推出关系 \(A=S\oplus r\)。

推论 1（韦德伯恩定理）. —— 设 K 是一个交换域，A 是一个 K-代数，r 是 A 的根。假设 K-代数 A/r 是绝对半单的。

\begin{enumerate}
\def\labelenumi{\alph{enumi})}
\item
  设 \(S_1\) 与 \(S_2\) 是 A 的满足 \(A = S_1 \oplus \mathfrak{r} = S_2 \oplus \mathfrak{r}\) 的子代数。存在 A 的一个把 \(S_1\) 变到 \(S_2\) 的特殊自同构。
\item
  若 \(\mathfrak{r}\) 是幂零的，则存在一个子代数 \(S\)（\(A\) 的），满足 \(A = S \oplus \mathfrak{r}\)。
\end{enumerate}

这由命题 7 与 VIII, 第 238 页的定理 2 得出。

推论 2. —— 设 A 是完满域 K 上的一个交换的有限次数代数，r 是它的根。存在 A 的唯一的子代数 S，使得 \(A = S \oplus r\)。此外，S 同构于 K 的有限多个有限次数扩张的一个乘积。

K-代数 A/r 是半单的（VIII, 第 173 页，命题 1）且次数有限；由于域 K 是完满的，它是绝对半单的（VIII, 第 232 页，定理 1）。由于理想 r 是幂零的且 A 是交换的，S 的存在性与唯一性便由推论 1 得出。由于 S 是半单、交换且次数有限的，最后的断言是 VIII, 第 137 页的命题 3 的一个推论。

注. —— 1) A/r 绝对半单这一假设在推论 1 中是本质的（VIII, 第 246 页，习题 4）。

\begin{enumerate}
\def\labelenumi{\arabic{enumi})}
\setcounter{enumi}{1}
\tightlist
\item
  假设 A 是域 K 上的一个阿廷代数。若 A 是交换的，则我们可以证明（VIII, 第 180 页，习题 9）A 同构于一个代数乘积 \(A_{1} \times \cdots \times A_{n}\)，使得对每个 i，\(\mathrm{A}_{i}/\Re(\mathrm{A}_{i})\) 都是一个域。反之，若 A 不是交换的，则 A 可能不同构于这样的代数乘积 \(A_{1} \times \cdots \times A_{n}\)：对每个 i，\(\mathrm{A}_{i}/\Re(\mathrm{A}_{i})\) 都是单环（VIII, 第 247 页，习题 5）。
\end{enumerate}

